% §4 Path-Based Tree Encodings

The index-set framework of \cref{sec:index-set} is concise and
sufficient for the three classical encodings.  However, it has a
structural limitation: deriving valid index sets from tree structures
requires strong structural conditions that most trees do not satisfy
(we prove this precisely in \cref{sec:star-tree}).  We now present a
second parameterization that removes this restriction entirely.

\subsection{Labelled rooted trees}

\begin{definition}[Encoding tree]
\label{def:encoding-tree}
An \emph{encoding tree} for $n$ fermionic modes is a rooted tree $T$
with internal nodes, where each node $v$ has at most three descending
connections.  Each connection is either an \emph{edge} (to a child
node) or a \emph{leg} (a dangling link to no child).  Every node has
exactly three connections, each labelled by a Pauli operator:
one $X$, one $Y$, one $Z$.  The total number of legs is $2n$ (one pair
per fermionic mode).
\end{definition}

\noindent
The labelling scheme allocates legs to modes: if a node has $k$ children
($0 \le k \le 3$), it contributes $3 - k$ legs.  The legs across the
entire tree are paired into $n$ mode assignments, with each mode
receiving one ``$x$-type'' and one ``$y$-type'' leg.

\subsection{Construction: from tree to encoding}

Given an encoding tree $T$, the encoding of mode $j$ proceeds as
follows.

\paragraph{Step 1: Identify legs.}
For mode $j$, let $\ell_j^{(x)}$ and $\ell_j^{(y)}$ be the paired
legs.

\paragraph{Step 2: Trace root-to-leg paths.}
For each leg, walk from the root to the leg, recording the Pauli
label of each edge traversed and appending the leg's own label at the
end.  Let $\pi(\ell)$ denote this sequence of Pauli labels.

\paragraph{Step 3: Construct Majorana strings.}
The Majorana operator for leg $\ell$ is the Pauli string obtained by
placing the path labels at the corresponding qubit positions (one
qubit per internal node) and identity elsewhere:
\begin{equation}
  S(\ell) = \bigotimes_{v \in \mathrm{path}(\ell)}
    \text{label}(v \to \text{child}) \;\;\otimes\;\;
    \bigotimes_{v \notin \mathrm{path}(\ell)} I.
  \label{eq:path-pauli}
\end{equation}

\paragraph{Step 4: Combine into ladder operators.}
\begin{equation}
  \ad{j} = \tfrac{1}{2}\bigl(S(\ell_j^{(x)}) - i\, S(\ell_j^{(y)})\bigr).
  \label{eq:tree-ladder}
\end{equation}

\begin{theorem}[Tree encoding correctness]
\label{thm:tree-car}
For any encoding tree $T$ with $2n$ legs paired into $n$ modes, the
ladder operators defined by~\eqref{eq:tree-ladder} satisfy the canonical
anticommutation relations~\eqref{eq:car1}--\eqref{eq:car3}.
\end{theorem}

\begin{proof}[Proof sketch]
The key property is that legs on distinct subtrees produce Pauli strings
with disjoint non-identity support except at the branching ancestor,
where the distinct labels ($X$, $Y$, or $Z$) ensure anticommutation.
Legs on the same subtree inherit a shared prefix and differ only in the
subtree below the branch point.  The $\Zi$ phase arithmetic of
\cref{sec:exactness} ensures that all cross-terms cancel exactly.  A
detailed proof by induction on tree depth follows the
construction of Jiang et al.~\cite{jiang2020} and Miller et
al.~\cite{miller2023bonsai}.
\end{proof}

\subsection{Recovery of known encodings}

Different tree topologies recover all five known encodings:

\begin{table}[h]
\centering
\caption{Tree topologies and the encodings they induce.}
\label{tab:tree-recovery}
\begin{tabular}{@{}lll@{}}
\toprule
Tree topology & Encoding recovered & Max Pauli weight \\
\midrule
Linear chain (star, depth 1) & Jordan--Wigner & $O(n)$ \\
Star with different ordering & Parity & $O(n)$ \\
Fenwick tree (depth $\lceil\log_2 n\rceil$) & Bravyi--Kitaev & $O(\log_2 n)$ \\
Balanced binary (depth $\lceil\log_2 n\rceil$) & Binary tree & $O(\log_2 n)$ \\
Balanced ternary (depth $\lceil\log_3 n\rceil$) & Ternary tree & $O(\log_3 n)$ \\
\bottomrule
\end{tabular}
\end{table}

\subsection{Pauli-weight bounds}

\begin{theorem}[Weight bound]
\label{thm:weight-bound}
For an encoding tree $T$ with depth $d = \depth{T}$, every encoded
operator $\ad{j}$ or $\an{j}$ has Pauli weight at most $2d + 1$.
\end{theorem}

\begin{proof}
Each ladder operator is a sum of two Pauli strings (from the two
Majorana operators).  Each Majorana string has weight equal to the
length of the root-to-leg path, which is at most $d + 1$ (depth plus
the leg label).  The worst case for the sum
$\frac{1}{2}(S^{(c)} \pm i S^{(d)})$ occurs when the two paths share
only the root, giving weight at most $2d + 1$.
\end{proof}

\begin{corollary}[Optimal asymptotic weight]
\label{cor:optimal-weight}
Among trees on $n$ modes:
\begin{enumerate}
  \item[(a)] The balanced binary tree achieves weight $O(\log_2 n)$.
  \item[(b)] The balanced ternary tree achieves weight $O(\log_3 n)$,
    which is optimal among fixed-arity trees since each node can have
    at most 3 children (one per Pauli label).
\end{enumerate}
\end{corollary}

\subsection{Relationship between the two frameworks}

The index-set and tree frameworks are complementary:

\begin{enumerate}
  \item \textbf{Index-set $\to$ tree.}  Given a valid scheme
    $\scheme$, one can construct a tree whose path-based encoding
    reproduces the same Majorana operators.  However, this requires
    the index sets to be compatible with a tree structure (not all
    set triples arise from trees).

  \item \textbf{Tree $\to$ index-set.}  Given a tree $T$, one can
    attempt to derive index sets $U$, $P$, $\mathrm{Occ}$ by
    reading off ancestor, child, and descendant relationships.
    This derivation (``Construction~A'' in the terminology
    of~\cite{seeley2012}) satisfies the CAR only when the tree is a
    \emph{star} (depth~$\le 1$), as we prove in \cref{sec:star-tree}.
    This is a far more restrictive domain than previously recognized.

  \item \textbf{Path-based is universal.}  The path-based construction
    produces valid encodings for \emph{every} rooted tree, regardless of
    monotonicity.  It is therefore the more general framework, while the
    index-set framework provides a compact algebraic characterization
    for the three classical encodings.
\end{enumerate}

\begin{remark}[The encoding space]
By Cayley's formula~\cite{cayley1889}, there are $n^{n-1}$ labelled
rooted trees on $n$ nodes, each defining a valid encoding via the
path-based construction.  The classical encodings are a measure-zero
subset of this exponentially large space.  Different trees in this space
trade off maximum Pauli weight, mean Pauli weight, and structural
properties such as symmetry and stabilizer structure.
\end{remark}
